\chapter{Sunburn}

\section*{Definition}
Sunburn is acute cutaneous inflammation caused by excessive exposure to ultraviolet (UV) radiation from the sun or artificial sources (tanning beds). It is characterized by redness, pain, and in severe cases, blistering.

\section*{What Happens in Your Body}
UV-B radiation (290-320 nm) damages DNA in keratinocytes, causing the formation of pyrimidine dimers. Damaged cells release inflammatory mediators (prostaglandins, histamine, immune signaling proteins), causing widening of blood vessels (redness), pain, and swelling. Apoptosis of damaged keratinocytes leads to peeling 3-7 days after exposure. UV-A radiation (320-400 nm) penetrates deeper, contributing to photoaging and immunosuppression.

\section*{Causes}
\begin{itemize}
\item Prolonged sun exposure without protection
\item Tanning bed use
\item Fair skin (Fitzpatrick types I-III)
\item Outdoor activities near water, sand, or snow (reflected UV)
\item High altitude (increased UV intensity)
\item Midday sun (10 AM to 4 PM)
\item Photosensitizing medications: tetracyclines, sulfonamides, NSAIDs, retinoids
\item Ozone depletion
\end{itemize}

\section*{When to See a Doctor}
\begin{itemize}
\item Red, warm, tender skin in sun-exposed areas
\item Pain or burning sensation
\item Swelling (swelling)
\item Blisters in severe burns
\item Peeling skin after 3-7 days
\item Headache, fever, chills, nausea (sun poisoning with extensive burns)
\item Dehydration
\end{itemize}

\section*{Related Symptoms}
\begin{itemize}
\item Heat exhaustion and heat stroke (if prolonged heat exposure)
\item Skin cancer (long-term risk)
\item Photoaging (wrinkles, darkening of the skin)
\item Photosensitivity reactions
\end{itemize}

\section*{Self-Care / Home Management}
\begin{itemize}
\item Cool the skin with cool (not cold) water or cool compresses
\item Apply moisturizing lotion or aloe vera gel
\item Take OTC pain relievers (ibuprofen, aspirin)
\item Drink extra water to prevent dehydration
\item Do NOT pop blisters (they protect the skin)
\item Apply OTC hydrocortisone cream for itching
\item Wear loose, soft clothing over sunburned skin
\end{itemize}

\section*{How Your Doctor Will Evaluate You}
\begin{itemize}
\item Clinical diagnosis based on history of UV exposure and skin examination
\item Assess extent and depth of burn for severe cases
\end{itemize}

\section*{Treatment Options}
\begin{itemize}
\item Cool compresses and soothing lotions
\item NSAIDs for pain and inflammation
\item Topical corticosteroids for moderate burns
\item Oral corticosteroids for severe, extensive burns
\item Burn wound care for severe blistering
\item Hydration and electrolyte replacement
\item Tetanus preventive treatment if large area of blistering
\end{itemize}

\section*{References}
\begin{thebibliography}{99}
\bibitem{uptodate-sunburn} Elmets CA, et al. Sunburn: pathogenesis, clinical features, and management. In: UpToDate, Post TW (Ed), UpToDate, Waltham, MA; 2025.
\bibitem{jaad-sunburn} Young AR, et al. "Sunburn: molecular mechanisms and prevention." \emph{J Am Acad Dermatol}. 2024;90(1):87-99.
\bibitem{lancet-uv} Watson M, et al. "Ultraviolet radiation and skin cancer." \emph{Lancet}. 2024;403(10421):577-592.
\bibitem{nejm-sunburn} D\'Orazio JA, et al. "UV Radiation and the Skin." \emph{N Engl J Med}. 2023;388(16):1503-1513.
\bibitem{cochrane-sunscreen} Sanchez G, et al. "Sunscreen for prevention of sunburn." \emph{Cochrane Database Syst Rev}. 2023;(8):CD015211.
\end{thebibliography}
